\documentclass[11pt]{amsart}
\usepackage[utf8]{inputenc}


\RequirePackage[numbers]{natbib}
\usepackage[dvips]{epsfig}
\usepackage{graphicx}
\usepackage{latexsym}
\usepackage{cancel}
\usepackage{ulem}
\usepackage[colorlinks=true,citecolor=blue,urlcolor=blue]{hyperref}

%\usepackage[notref, notcite]{showkeys}
\usepackage{amsmath,amsthm,amsfonts,amssymb,amscd,amsbsy,dsfont,hyperref}

\usepackage[margin=3cm]{geometry}
\setlength{\oddsidemargin}{.5cm} 
\setlength{\evensidemargin}{.5cm}
\setlength{\textwidth}{15cm} 
\setlength{\textheight}{20cm}
\setlength{\topmargin}{1cm}

\usepackage{aquabox}
\usepackage{stmaryrd}
\usepackage{graphicx}
\usepackage{float}
\usepackage{centernot}

\makeatletter
\def\l@section{\@tocline{1}{12pt plus2pt}{0pt}{}{\bfseries}}% <- added

\def\@tocline#1#2#3#4#5#6#7{\relax
    \ifnum #1>-1% <- added
  \ifnum #1>\c@tocdepth % then omit
  \else
    \par \addpenalty\@secpenalty%
    \begingroup \hyphenpenalty\@M
    \@ifempty{#4}{%
      \@tempdima\csname r@tocindent\number#1\endcsname\relax
    }{%
      \@tempdima#4\relax
    }%
    \parindent\z@ \leftskip#3\relax \advance\leftskip\@tempdima\relax
    \rightskip\@pnumwidth plus4em \parfillskip-\@pnumwidth
    #5\leavevmode\hskip-\@tempdima #6\nobreak\relax
    \hfil\hbox to\@pnumwidth{\@tocpagenum{#7}}\par
    \nobreak
    \endgroup
  \fi
\fi}
\makeatother
\setcounter{tocdepth}{2}  



\usepackage{enumitem}
\newcommand{\subscript}[2]{$#1 _ #2$}
\newcommand{\delete}[1]{}

\newcommand{\E}{\espe}
\newcommand{\R}{\RR}
\renewcommand{\Re}{\mathfrak{R}e}

\newcommand{\norm}[1]{\lVert#1\rVert_2}

\newcommand{\rj}{\ln^2(J)}
\newcommand{\redline}{\noindent\textcolor{red}{\rule{16cm}{1mm}}}
\newcommand{\ul}[1]{\underline{#1}}

\newcommand{\change}[1]{\textcolor{red}{#1}}

\newcommand{\mmax}{m_{\max}}
\newcommand{\tdiff}{(t_2 - t_1)}
\newcommand{\mopt}{m_J^\star}

\newcommand{\diff}{\,\mathrm{d}}

\newenvironment{details}
  {\cor{violet}}%
  {}
\newenvironment{review}%
  {\color{red}}%
  {}
 \newenvironment{review1}%
  {\color{blue}}%
  {}

\newcommand{\supp}{\mathrm{supp}\,}
\newcommand{\spec}{\mathrm{spec}\,}
\newcommand{\tendsto}[1]{\xrightarrow[#1]{}}

\title{A mutation equation}
\author{}
\date{\today}

\begin{document}

\maketitle

\tableofcontents

% \section{The Mutation Equation : all mutation are deleterious}

% \subsection{Description of the model}

% Consider the \textit{mutation equation}
% \begin{equation}\left\{
%     \begin{array}{ll}
% 	\frac{\partial}{\partial t}u(t,x) + \lambda u(t,x) = \lambda \int_{0}^{\infty} k(y, x) u(t,y) \dy
% 	\eqsepv
% 	& t \in (0, \infty), x \in \RR_+, \lambda \in (0, \infty) \eqfinv
% 	\\
% 	u(0, x) = u_0(x)
% 	& x \in \RR_+
% 	\eqfinp
%     \end{array}
% 	\right.
% 	\label{eq:pseudo}
% \end{equation}

% In all that follows, we will assume that $k$ is self-similar, \ie 
%     \begin{equation*}
%         k(y, x) = \frac{1}{y} k_0 \Big(\frac{x}{y} \Big)
%         \eqsepv
%         y > 0
%         \eqsepv 
%         x \in \RR_+
%         \eqfinp
%     \end{equation*}
%     \begin{equation*}
%         \supp k_0 \subset [0, 1] 
%         \eqsepv
%         \boxed{\int_0^1 k_0(z) \dz = 1}
%         \eqfinp
%     \end{equation*}

% In addition, we assume that 
%     \begin{equation}
%         \boxed{\int_0^1 k_0(z) (1 + z) \dz < \infty}
%         \eqfinp
%         \label{H:douesco4}
%     \end{equation}
% \onlyME{En cas de besoin, on supposera l'existence de moments pour tout $n \in \NN$.}
    
% \begin{remark}
%     This last assumption ensures that the Mellin transform $K(s) := \Mcal_{k_0}(s) = \int_0^\infty k_0(z)z^{s-1}\dz$ is well defined for all $\Re(s) \in (1, 2)$.
% \end{remark}
    
% Under theses assumptions, Equation\eqref{eq:pseudo} becomes
%     \begin{equation}
%     	\frac{\partial}{\partial t}u(t,x) + \lambda u(t,x) = \lambda \int_{x}^{\infty} \frac{1}{y} k_0 \Big(\frac{x}{y} \Big) u(t,y) \dy
% 	    \eqfinp 
% 	    \label{eq:similar_pseudo_growth_equation}
%     \end{equation}

% \subsubsection{Explicit solution to the mutation Equation \eqref{eq:similar_pseudo_growth_equation}}

% We assume that     \begin{equation}
%         u_0 \geq 0 \mtext{and} \int_0^\infty u_0(x) (1+x) \dx < \infty
%         \eqfinp
%         \label{H:u0_L1}
%     \end{equation}
    
% Define 
%     \begin{equation}
%         I(u_0) := \left\{ p \in \RR \eqsepv \int_0^\infty u_0(x) x^{p-1} \dx < \infty \right\}
%         \eqfinp
%     \end{equation}
% By Assumption~\eqref{H:u0_L1}, $[1, 2] \subset I(u_0)$. Moreover $I(u_0)$ is a interval of $\RR$.\\
% It leads us to define
%     \begin{subequations}
%     \begin{align*}
%         p_0 &= \inf \{p : p \in I(u_0) \} \in [-\infty, 1] \eqfinv\\
%         q_0 &= \sup \{q : q \in I(u_0) \} \in [2, \infty] \eqfinp
%     \end{align*}    
%     \end{subequations}

% For the simplicity, we assume that 
%     $ p_0=-\infty$ and $q_0 = \infty$
    
% We also assume that 
%     \begin{equation}
%         \forall \nu \in (p_0, q_0) 
%         \eqsepv
%         \int_0^\infty |u_0'(x)| x^\nu \dx < \infty
%         \eqsepv
%         \lim_{x \to 0} x^\nu u_0(x) = 0
%         \eqsepv
%         \lim_{x \to \infty} x^\nu u_0(x) = 0
%         \label{H:DoumicEsco22}
%     \end{equation}
%     \begin{equation}
%         \forall \nu \in (p_0, q_0) 
%         \eqsepv
%         \int_0^\infty |u_0''(x)| x^\nu \dx < \infty
%         \eqsepv
%         \lim_{x \to 0} x^{\nu + 1} u_0'(x) = 0
%         \eqsepv
%         \lim_{x \to \infty} x^{\nu + 1} u_0'(x) = 0
%         \label{H:DoumicEsco23}
%     \end{equation}
    
% In the following, we consider the Mellin transforms
%     \begin{equation*}
%         U(t,s) := \Mcal_{u(t, \cdot)}(s)
%         \mtext{and}
%         K(t,s) := \Mcal_{k_0(t, \cdot)}(s)
%     \end{equation*}
% when these integrals are well-defined.

% \begin{theorem} Let 
% $u(t, x) \in \Ccal([0, \infty); L^1((1+x)\dx)) \cap \Ccal^1([0, \infty); L^1((1+x)\dx))$ be a solution of Equation\eqref{eq:similar_pseudo_growth_equation}. Then 
% 	\begin{equation*}
% 		\frac{\partial}{\partial t} U(t,s) + \lambda U(t,s) = \lambda K(s) U(t,s)
% 		\eqfinp
% 	\end{equation*}
% \end{theorem}
% \begin{proof}
% Let multiply Equation\eqref{eq:similar_pseudo_growth_equation} by $x^{s-1}$ and integrate over $(0, \infty)$, we obtain
% 	    \begin{align*}
%     	   \frac{d}{dt} \int_0^\infty u(t,x) x^{s-1} \dx
%     	   + 
%     	   \lambda \int_0^\infty u(t,x)x^{s-1} \dx
%     	   &= 
%     	   \lambda \int_0^\infty x^{s-1} \int_{x}^{\infty} \frac{1}{y} k_0 \Big(\frac{x}{y} \Big) u(t,y) \dy \dx \\
% 	       &= 
% 	       \lambda \int_0^\infty  u(t,y) \int_{0}^{y} \frac{1}{y} k_0 \Big(\frac{x}{y} \Big) x^{s-1} \dx \dy 	   \\
% 	       &= 
% 	       \lambda \int_0^\infty  u(t,y) y^{s-1}  \int_{0}^{1} k_0(z) z^{s-1} \dz \dy	   \\
% 	       &= \lambda \int_0^\infty  u(t,y) y^{s-1} \dy  \int_{0}^{1} k_0(z) z^{s-1} \dz
% 	       \eqfinp
% 	    \end{align*}
% Then 
% 	\begin{equation*}
% 		\frac{\partial}{\partial t} U(t,s) + \lambda U(t,s) = \lambda K(s) U(t,s)
% 		\eqfinp
% 	\end{equation*}
	
% Then 
%     \begin{equation*}
%         U(t,s) = U_0(s) e^{\lambda(K(s) - 1) t}
%         \eqsepv
%         U_0(s) = \Mcal_{u_0}(s)
%         \eqfinp
%     \end{equation*}
% \end{proof}

% By using the inverse Mellin transform, we have at least formally the explicit solution

% \begin{equation}
%     u(t,x) = \frac{1}{2i\pi} \int_{c - i \infty}^{c + i \infty} U_0(s) e^{\lambda(K(s) - 1) t} x^{-s}\ds 
%     \eqsepv x \in (0, \infty) \as
%     \label{eq:explicit_reconstruction}
% \end{equation}

% \begin{theorem}(Existence of an explicit solution) Assume \eqref{H:u0_L1}, \eqref{H:DoumicEsco22} and \eqref{H:DoumicEsco23}.
% Then, there exists a unique function $u \in \Ccal([0, \infty); L^1((1+x)dx)) \cap \Ccal^1((0, \infty); L^1((1+x)dx))$ that satisfies Equation~\eqref{eq:similar_pseudo_growth_equation} for all $t > 0$ and almost every $x > 0$. This solution is given by \eqref{eq:explicit_reconstruction} for all $\nu \in (p_0, q_0)$.
% \end{theorem}
% \begin{proof}
% The result is a consequence of the Mellin inversion formula.
% In order to apply it, we have to prove that for any $\nu \in (1, 2)$, 
%     \begin{equation*}
%         \int\limits_{\nu - i \infty}^{\nu + i \infty} \bigg| U_0(s) e^{\lambda(K(s) - 1) t} x^{-s} \bigg| \ds < \infty
%         \eqfinp
%     \end{equation*}
% We know that for all $s \in [1, 2]$, 
%     \begin{align*}
%         K(s) = \int_{0}^{\infty} k_0(z) z^{s-1} \dz
%         &= \int_{0}^{1} k_0(z) z^{s-1} \dz \\
%         &\leq \int_{0}^{1} k_0(z) \dz
%         \\
%         &= K(1)  \\
%         &= 1
%         \eqfinp
%     \end{align*}
    
% Then 
%     \begin{equation*}
%          \bigg| U_0(s) e^{\lambda(K(s) - 1) t} x^{-s} \bigg|  <   \bigg| U_0(s)  x^{-\nu} \bigg|
%     \end{equation*}

% It only remains to prove that 
%     \begin{equation}
%             \int\limits_{\nu - i \infty}^{\nu + i \infty} \big| U_0(s)\big| \ds 
%             = 
%             \int\limits_{-\infty}^{\infty} \big| U_0(\nu + is)\big| \ds
%             < \infty
%             \label{eq:DouEsc43}
%     \end{equation}

% This quantity has already been studied under the same assumptions by Doumic and Escobedo \cite{doumic2016time}. Integration by parts lead to 
%     \begin{align*}
%         U_0(\nu + i v) 
%         &= \int_0^\infty u_0(x) x^{\nu-1}x^{i v}\dx \\
%         &= \frac{1}{iv + 1} \big[ u_0(x) x^{\nu} x^{i v} \big]_0^\infty - \frac{1}{iv + 1} \int_0^\infty (u_0(x) x^{\nu-1})' x^{i v+1}\dx \\
%         \eqfinp
%     \end{align*}
% Assumption \eqref{H:DoumicEsco22} guarantees that $[ u_0(x) x^{\nu} x^{i v} ]_0^\infty = 0$. Therefore
%     \begin{align*}
%         U_0(\nu + i v) 
%         &= - \frac{1}{iv + 1} \int_0^\infty (u_0(x) x^{\nu-1})' x^{i v+1}\dx \\
%         &= - \frac{\nu - 1}{iv + 1} \int_0^\infty u_0(x) x^{\nu-2} x^{i v + 1}\dx  - \frac{1}{iv + 1} \int_0^\infty u_0'(x) x^{\nu-1} x^{i v + 1}\dx \\
%         &= - \frac{\nu - 1}{iv + 1} \int_0^\infty u_0(x) x^{\nu-1} x^{i v}\dx  - \frac{1}{iv + 1} \int_0^\infty u_0'(x) x^{\nu} x^{i v}\dx \\
%         &= - \frac{\nu - 1}{iv + 1} U_0(\nu + i v)   - \frac{1}{iv + 1} \int_0^\infty u_0'(x) x^{\nu} x^{i v}\dx
%         \eqfinp
%     \end{align*}
% It follows that 
%     \begin{align*}
%         U_0(\nu + i v)
%         &= - \frac{1}{(iv + \nu)}\int_0^\infty u_0'(x) x^{\nu} x^{i v}\dx
%         \eqfinp
%     \end{align*}
% \onlyME{It follow that 
%     \begin{equation*}
%         |U_0(\nu + i v)| \leq \frac{1}{\sqrt{\nu^2 + v^2}} \int_0^\infty |u_0'(x)| x^{\nu} \dx
%         \eqfinp
%     \end{equation*}}
% The same calculations and Assumption\eqref{H:DoumicEsco23} allow us to obtain 
%     \begin{align*}
%         U_0(\nu + i v)
%         &=  \frac{1}{(iv + \nu)(iv + \nu + 1)}\int_0^\infty u_0''(x) x^{\nu + 1} x^{i v}\dx
%         \eqfinp
%     \end{align*}
% It follow that 
%     \begin{equation*}
%         |U_0(\nu + i v)| \leq \frac{1}{\sqrt{(\nu^2 + v^2)((1+\nu)^2 + v^2)}} \int_0^\infty |u_0''(x)| x^{\nu + 1} \dx
%         \eqfinp
%     \end{equation*}

% We have that $\int_0^\infty |u_0''(x)| x^{\nu + 1} \dx < \infty$ from Assumption\eqref{H:DoumicEsco23}. It allows us to deduce \eqref{eq:DouEsc43} \\

% Conversely, if $u_0$ satisfies \eqref{H:u0_L1}, \eqref{H:DoumicEsco22} and \eqref{H:DoumicEsco23}, then 
%  $$u \in \Ccal([0, \infty); L^1((1+x)dx)) \cap \Ccal^1((0, \infty); L^1((1+x)dx))$$
% where $u$ is defined by \eqref{eq:explicit_reconstruction}. By construction, it is a solution of the equation \eqref{eq:similar_pseudo_growth_equation}.
% \end{proof}

% \subsection{Speed of convergence}

% In the following, we define
%     \begin{equation*}
%         w(t,x) := e^t u(t,x)
%     \end{equation*}
% Then $\forall c \in (1, 2)$
%     \begin{equation}
%         w(t,x) = \frac{1}{2i\pi} \int_{c - i \infty}^{c + i \infty} U_0(s) e^{\lambda K(s) t} x^{-s}\ds 
%     \eqsepv x \in (0, \infty) \as
%     \end{equation}

% We define $s_+(t,x) := (K')^{-1} \Big( \frac{\log x}{\lambda t}\Big)$

% \begin{theorem}
%     Assume that $u_0$ satisfies the condition \ref{H:u0_L1}, \ref{H:DoumicEsco22} and \ref{H:DoumicEsco23} and let $\varepsilon(t)$ such that $\varepsilon(t)\to 0$ but $t\varepsilon^2(t) \to \infty$ as $t\to \infty$\\

%  Suppose that  the measure $k_0$ has a non zero absolutely continuous part. Then for all $\delta >0$ arbitrarily small, there exists two functions $\gamma _{ \delta  } (\varepsilon )$ and $\omega _{ \delta  } (\varepsilon )$ such that:

% \begin{eqnarray}
% &&u(t, x)=x^{-s_+(t, x)}e^{(K(s_+(t, x))-1)t}\times \nonumber\\
% &&\hskip 3cm \left(  \frac {U_0(s_+(t, x))+\omega _{ \delta  } (\varepsilon )} {\sqrt{2\pi t K''(s_+)}}\,\Theta_1(t)+\mathcal O\left(e^{- t\gamma  _{ \delta  }(\varepsilon (t)) }\right)  \right) \label{theorem3:1}\\
% &&\hskip 1.5cm \hbox{as}\,\,t\to \infty,\,\,\hbox{uniformly for}\; x\;\hbox{such that}\,\,  p_0+\delta<s _+(t, x)<q_0-\delta, \nonumber\\
% &&\Theta_1(t)=\left(1+
%  \mathcal O\left(\frac {e^{-\frac {t\varepsilon^2 (t)} {2}K''(q_0)}} {\sqrt {t \varepsilon^2 (t)K''(q_0)}}\right)\right),\,\,\,\hbox{as}\,\,t\to \infty\label{theorem3:2}
% \end{eqnarray}
% where,  $\lim _{ \varepsilon \to 0}\varepsilon  ^{-2}\gamma _{ \delta  }  (\varepsilon  )=C _{ \delta  }$ for some constant $C _{ \delta  }\in (K^{''} (q_0) /2 ,  K^{''} (p_0)/2)$. The function  $\omega _{ \delta  }(\varepsilon )$, defined later, satisfies
% \begin{eqnarray}
% \label{theorem3:omega}
% &&|\omega _{ \delta  } (\varepsilon )|\le C (p_0+\delta,q_0-\delta)\left( \int _{R(\varepsilon )}^\infty \left|z^{q_0-\delta-1}u_0(z)\right|dz+\right. \nonumber \\
% &&\hskip 4cm \left. +\int _0^{\frac {1} {R(\varepsilon )}} \left|z^{p_0+\delta-1}u_0(z)\right|dz+\varepsilon  \log R(\varepsilon )\right),
% \end{eqnarray}
% where the constant $C(p_0+\delta,q_0+\delta)$ depends on the integrability properties of $u_0$, $R$ is any function such that $R(\varepsilon )\to \infty$, $\varepsilon \log R(\varepsilon )\to 0$ as $\varepsilon \to 0$.
% \end{theorem}

% \subsubsection{Proof}

% Here, we want to apply the \textit{steepest descent method} \cite{debye1909naherungsformeln} when $t \to \infty$.

% We recall that 
%     \begin{equation*}
%         \omega(t,x)
%         = \int_{c-i\infty}^{c+i\infty} U_0(s) e^{\lambda K(s)t}x^{-s} \ds
%         =  \int_{c-i\infty}^{c+i\infty} U_0(s) e^{\phi(s, t, x)} \ds
%         \eqfinp
%     \end{equation*}
% with $\phi(s,t,x) := \lambda K(s) t - s \ln(x) $.\\

% In order to apply the steepest descent method, our first step consist in finding the stationary point of $\phi'(\cdot, t, x) = 0$. 

% Define $p_0 < s_+(t,x) := (K')^{-1} \Big( \frac{\log x}{\lambda t} \Big) < q_0$.

% \paragraph{\ul{Step 1 -- Show that $s \mapsto \Re(\phi(s, t, x))$ reachs its maximum at $s_+(t, x)$}} \text{}\\

% $\blacktriangleright$ Consider a neighborhood of $s_+(t,x)$ : $|s - s_+(t,x)| < \varepsilon$. We have
% \begin{equation}
%     K(s) = K(s_+) + (s - s_+(t,x)) K'(s_+(t,x)) + \frac{(s - s_+(t,x))^2}{2} K''(s_+(t,x)) + \frac{(s - s_+(t,x))^3}{6} K'''(c(t,x))
% \end{equation}
% where $s < c(t,x) < s_+(t,x)$.
%     \begin{align*}
%          |K'''(c)|
%          &= |\int_0^1 (\log z)^3 k_0(z) e^{(c-1)\log(z)}\dz|
%          \leq \int_0^1 |(\log z)^3| k_0(z)  e^{\Re (c-1)\log(z)}\dz \\
%          &\leq  \boxed{\int_0^1 |(\log z)^3| k_0(z)  e^{(q_0-1)\log(z)}\dz < \infty}
%     \end{align*}
    
% Given that $|s - s_+| < \varepsilon$ and $s_+ + \varepsilon < q_0$, we have 
%     \begin{equation}
%         \Re K(s) = K(s_+) + \frac{(s - s_+(t,x))^2}{2} K''(s_+(t,x)) + O(\varepsilon^3)
%         \eqfinp
%     \end{equation}
% and 
%     \begin{equation}
%         \Re \phi(s, t, x) = \phi(s_+, t, x) + \lambda t \Big(\frac{(s - s_+(t,x))^2}{2} K''(s_+(t,x)) + O(\varepsilon^3)\Big)
%         \eqfinp
%     \end{equation}

    
% $\blacktriangleright$ Outside the neighborhood of $s_+(t,x)$ : $|s - s_+(t,x)| < \varepsilon$. We have
%     \begin{align*}
%         \Re(\phi(s_+ + i \nu, t, x) - \phi(s_+, t, x)) 
%         &=
%         \Re \bigg( \int_0^1 x^{s_+ - 1}k_0(x) (x^{i\nu} - 1) \dx \bigg) \\
%         &=
%         \Re \bigg( \int_0^1 x^{s_+ - 1}k_0(x) (e^{i\nu \log(x)}- 1) \dx \bigg) \\ 
%         &= 
%         \int_0^1 x^{s_+ - 1}k_0(x) [\cos(\nu \log(x))- 1] \dx  \\
%         &=: G(v) \leq 0
%         \eqfinp
%     \end{align*}
    
% Recall the that $k_0(x) = g(x)dx + \dmu + \dnu$ where
% $g \in L^1(\dx)$
% $\dmu$ is a singular continuous measure and $\dnu$ a singular discrete measure.\\ 

% In the following, we assume that $k_0(x) :=  g(x)dx$.\\

% For any $\nu \in \RR \setminus \{0\}$, we have
% \begin{equation}
%      \int_0^1 x^{s_+ - 1}g(x) [\cos(\nu \log(x))] \dx < \int_0^1 x^{s_+ - 1}g(x) \dx
%      \eqfinp
% \end{equation}

% It follows that $G(\nu) < 0$, then there exists $\varepsilon > 0$, $R >0$ and $\delta > 0$
%  such that $\sup_{\varepsilon < \nu < R}G(\nu) < - \delta$.\\
 
%  It remains to verify that $\lim_{\nu \to +\infty}G(\nu) \neq 0$. 
%     \begin{align*}
%         G(\nu) 
%         &=  \int_0^1 x^{s_+ - 1}k_0(x) [\cos(\nu \log(x))- 1] \dx \\
%         &=  \int_{0}^\infty e^{-s_+ y}k_0(e^{-y})(\cos(-\nu y) - 1 )\dy
%         \end{align*}
% From Riemann–Lebesgue lemma, we have that 
%     \begin{equation}
%         \lim_{\nu \to +\infty}G(\nu) = - \int_{0}^\infty e^{-s_+ y}k_0(e^{-y})\dy = - \int_{0}^1 x^{s_+ - 1}k_0(x)\dx < 0
%     \end{equation}
    
% Therefore
% \begin{lemma} For all $\varepsilon > 0$, there exist $\gamma(\varepsilon)$ such
%     \begin{equation}
%         \sup_{|\nu| \geq \varepsilon } \Re(\phi(s_+ + i \nu, t, x)) \leq \Re(\phi(s_+, t, x)) - \gamma(\varepsilon)
%         \eqfinp
%     \end{equation}
% and
%     \begin{equation}
%         \sup_{|\nu| \leq \varepsilon } \Re(\phi(s_+ + i \nu, t, x)) \leq \Re(\phi(s_+, t, x))
%         \eqfinp
%     \end{equation}
% \end{lemma}

% \paragraph{\ul{Step 2 -- Use of the right path}} \text{}\\

% Close to $s_+(t,x)$, we consider the quantity
%     \begin{align*}
%     \frac{1}{2i \pi} \int_{s_+(t,x) - i\varepsilon}^{s_+(t,x) + i\varepsilon} U_0(s) e^{\phi(s,t,x)} \ds 
%     &=
%     \frac{1}{2i \pi} U_0(s_+(t,x)) \int_{s_+(t,x) - i\varepsilon}^{s_+(t,x) + i\varepsilon} e^{\phi(s,t,x)} \ds
%     \\
%     & \qquad \qquad \qquad+ 
%     \frac{1}{2i \pi} \int_{s_+(t,x) - i\varepsilon}^{s_+(t,x) + i\varepsilon} (U_0(s) -  U_0(s_+(t,x)))  e^{\phi(s,t,x)} \ds \\
%     &= I_1 + I_2
%     \end{align*}
    
% $\blacktriangleright$ \ul{The terme $I_1$}
% For any $s \in (s_+ - i \infty, s_+ + i \infty)$ such that $|s - s_+| < \varepsilon$
% We know that 
%     \begin{equation}
%         \Re \phi(s, t, x) = \phi(s_+, t, x) + t \Big(\frac{(s - s_+(t,x))^2}{2} K''(s_+(t,x)) + O(\varepsilon^3)\Big)
%         \eqfinp
%     \end{equation}

% Then 
% \begin{align*}
%    \int_{s_+(t,x) - i\varepsilon}^{s_+(t,x) + i\varepsilon} e^{\Re \phi(s,t,x)} \ds
%     &=
%     e^{\phi(s_+, t, x)} \cdot \int_{s_+(t,x) - i\varepsilon}^{s_+(t,x) + i\varepsilon} e^{t \Big(\frac{(s - s_+(t,x))^2}{2} K''(s_+(t,x)) + O(\varepsilon^3)\Big)} \ds \\
%     &= 
%      i e^{\phi(s_+, t, x)} \cdot \int_{- \varepsilon}^{\varepsilon} e^{t \Big(\frac{(s - s_+(t,x))^2}{2} K''(s_+(t,x)) + O(\varepsilon^3)\Big)} \ds \\
%     &= 
%      i e^{\phi(s_+, t, x)} \cdot \int_{- \varepsilon}^{\varepsilon} e^{- t \Big(\frac{v^2}{2} K''(s_+(t,x)) + O(\varepsilon^3)\Big)} \ds 
% \end{align*}

% Now we can see that we are back to the classical method of Laplace. The main idea is to transform the previous integral into a Gauss integral to have
%     \begin{equation}
%     \int_{- \varepsilon}^{\varepsilon} e^{- t \Big(\frac{v^2}{2} K''(s_+(t,x)) + O(\varepsilon^3)\Big)} \ds 
%     \sim_{t \to \infty}
%     \sqrt{2\pi} \frac{1}{\sqrt{tK''(s_+)}}
%     \end{equation}
    
% Indeed, 
% \begin{align*}
%     \int_{- \varepsilon}^{\varepsilon} e^{- t \Big(\frac{v^2}{2} K''(s_+(t,x)) + O(\varepsilon^3)\Big)} 
%     &= 
%      \frac{1}{\sqrt{tK''(s_+)}} \int_{- \varepsilon \sqrt{tK''(s_+) + \Ocal(\varepsilon})}^{\varepsilon \sqrt{tK''(s_+) + \Ocal(\varepsilon})} e^{- \frac{y^2}{2}}\dy \\
%      \intertext{and}
%      \int_{- \varepsilon \sqrt{tK''(s_+) + \Ocal(\varepsilon})}^{\varepsilon \sqrt{tK''(s_+) + \Ocal(\varepsilon})} e^{- \frac{y^2}{2}}\dy &= 
%     \sqrt{2\pi} - 2 e^{-\frac{t(\varepsilon^2 K''(s_+) + \Ocal(\varepsilon^3))}{2}} \cdot \Ocal\bigg(\frac{1 }{\varepsilon \sqrt{tK''(s_+) + \Ocal(\varepsilon})}\bigg)
%     \eqfinp
% \end{align*}
% It follows that 
%     \begin{align*}
%      \frac{1}{2i\pi}   \int_{s_+(t,x) - i\varepsilon}^{s_+(t,x) + i\varepsilon} e^{\Re \phi(s,t,x)} \ds 
%      &= 
%      \frac{1}{2i\pi} i e^{\phi(s_+, t, x)} \cdot \int_{- \varepsilon}^{\varepsilon} e^{- t \Big(\frac{v^2}{2} K''(s_+(t,x)) + O(\varepsilon^3)\Big)} \ds  \\
%      &= 
%       \frac{1}{2\pi} e^{\phi(s_+, t, x)} \cdot \frac{1}{\sqrt{tK''(s_+)}} \int_{- \varepsilon \sqrt{tK''(s_+) + \Ocal(\varepsilon})}^{\varepsilon \sqrt{tK''(s_+) + \Ocal(\varepsilon})} e^{- \frac{y^2}{2}}\dy \\
%      &=
%     \frac{1}{2\pi} e^{\phi(s_+, t, x)} \cdot \frac{1}{\sqrt{tK''(s_+)}} \bigg( 
%     \sqrt{2\pi} - 2 e^{-\frac{t(\varepsilon^2 K''(s_+) + \Ocal(\varepsilon^3))}{2}} \cdot \Ocal\bigg(\frac{1 }{\varepsilon \sqrt{tK''(s_+) + \Ocal(\varepsilon})}\bigg) \bigg) \\
%     \end{align*}
    
% Therefore, 
%     \begin{equation}
%     I_1 = \frac{U_0(s_+(t,x))}{\sqrt{2\pi}} e^{\phi(s_+, t, x)} \cdot \frac{1}{\sqrt{t \lambda K''(s_+)}} \bigg( 
%     1 + \Ocal\bigg(\frac{e^{-\frac{t(\varepsilon^2 \lambda K''(q_0))}{2}} }{\varepsilon \sqrt{\lambda tK''(q_0)}}\bigg) \bigg)\\
%     \end{equation}
    
% $\blacktriangleright$ \ul{The terme $I_2$}

% We have 
%     \begin{equation*}
%         I_2 = \frac{1}{2i \pi} \int_{s_+(t,x) - i\varepsilon}^{s_+(t,x) + i\varepsilon} (U_0(s) -  U_0(s_+(t,x)))  e^{\phi(s,t,x)} \ds 
%     \end{equation*}
    
% Here, we want to go back more or less to the previous case. 

%     \begin{align*}
%          |I_2| 
%          &= 
%          \frac{1}{2i \pi} \sup_{|v| \leq \varepsilon} \big[U_0(s_+ + iv) -  U_0(s_+) \big] \int_{s_+(t,x) - i\varepsilon}^{s_+(t,x) + i\varepsilon} \big| e^{\phi(s,t,x)} \big| \ds \\
%          &=
%          \sup_{|v| \leq \varepsilon} \big[U_0(s_+ + iv) -  U_0(s_+) \big] \frac{1}{\sqrt{2\pi}} e^{\phi(s_+, t, x)} \cdot \frac{1}{\sqrt{t \lambda K''(s_+)}} \bigg( 
%     1 + \Ocal\bigg(\frac{e^{-\frac{t(\varepsilon^2 \lambda K''(q_0))}{2}} }{\varepsilon \sqrt{\lambda tK''(q_0)}}\bigg) \bigg)
%     \end{align*}
    
% It only remains to study $  \sup_{|v| \leq \varepsilon} \big[U_0(s_+ + iv) -  U_0(s_+) \big]$ \\

% By Definition
% \begin{align*}
%     \big[U_0(s_+ + iv) -  U_0(s_+) \big]
%     &= \Big| \int_0^\infty  z^{s_+ - 1}u_0(z) (z^{iv} - 1) \dz \Big| \\
%     &\leq \int_0^{\frac{1}{R}} \Big|  z^{s_+ - 1}u_0(z) (z^{iv} - 1) \Big| \dz
%         + \int_{\frac{1}{R}}^R \Big|  z^{s_+ - 1}u_0(z) (z^{iv} - 1) \Big| \dz \\
%         &\qquad \qquad \qquad \qquad
%         + \int_R^{+\infty} \Big|  z^{s_+ - 1}u_0(z) (z^{iv} - 1) \Big| \dz \\
%     &\leq 2 \int_0^{\frac{1}{R}} \Big|  z^{p_0 + \delta - 1}u_0(z) \Big| \dz
%         +  \int_{\frac{1}{R}}^R \Big|  z^{s_+ - 1}u_0(z) (z^{iv} - 1) \Big| \dz \\
%         &\qquad \qquad \qquad \qquad
%         + 2 \int_R^{+\infty} \Big|  z^{q_0 - \delta - 1}u_0(z) \Big| \dz \\
% \end{align*}

% If $R \to + \infty$, then the first and the third terms vanish. Therefore, we only need to consider the integral $ \int_{\frac{1}{R}}^R \Big|  z^{s_+ - 1}u_0(z) (z^{iv} - 1) \Big| \dz$. We have
%     \begin{align*}
%          \int_{\frac{1}{R}}^R \Big|  z^{s_+ - 1}u_0(z) (z^{iv} - 1) \Big| \dz 
%          &\leq 
%           (\varepsilon \log R)\int_{\frac{1}{R}}^R \Big|  z^{s_+ - 1}u_0(z) \Big| \dz \\
%          &\leq 
%           (\varepsilon \log R)\int_{\frac{1}{R}}^R \Big|  (z^{p_0 + \delta - 1} + z^{q_0 - \delta - 1}) u_0(z) \Big| \dz \\
%     \end{align*}

% \subsection{Interpretation}


% We may rewrite 
% $$x u(t,x)\approx A(t,s_+) e^{t F(s_+(t,x))}$$
% where we define $F$ by
% \begin{equation*}
% F(s)={ K}( s)-1-(s-1) K'\left(s\right).
% \end{equation*}

% The behavior of $xu$ depends on the signe of $F$. 

\section{The Mutation Equation : mutation can be deleterious or beneficial}

\subsection{Location of the population}

Consider the equation
\begin{equation}\left\{
    \begin{array}{ll}
	\frac{\partial}{\partial t}u(t,x) = - \lambda u(t,x) + \lambda \int_{0}^{\infty} k\big( \frac{x}{y}\big)u(t,y) \frac{\dy}{y}
	\eqsepv
	& t \in (0, \infty), x \in \RR_+, \lambda \in (0, \infty) \eqfinv
	\\
	u(0, x) = u_0(x)
	& x \in \RR_+
	\eqfinp
    \end{array}
	\right.
    \label{eq:main}
\end{equation}

We assume that the initial data $u_0(x)$ is a function satisfying the following conditions 
$$ u_0 \geq 0 \eqsepv \int_0^\infty u_0(x)(1+x)\diff x < \infty.$$
It follows that the Mellin transform of $u_0$, $U_0(s)= \int_0^\infty x^{s-1}u(t,x)\diff x$ is analytic on the strip $\Re(\alpha) \in (p_0,q_0)$ where $p_0 \leq 1 < 2 \leq q_0$ \\
For the simplicity, we assume that 
$ p_0=-\infty$ and $q_0 = \infty$\\

By applying the Mellin transform for Equation~\eqref{eq:main}, we have
	\begin{equation}
		\frac{\partial}{\partial t} U(t,s) = - \lambda U(t,s) + \lambda K(s) U(t,s)
		\eqfinp
        \label{eq:explicit_reconstruction}
	\end{equation}
If we have some good hypothesis (see Section~\ref{s:explicit}) on $U_0$ and $K$, then we have an explicit formula for the inverse Mellin transform and it follows that
    \begin{equation*}
        U(t,s) = U_0(s) e^{-\lambda t(1 - K(s) ) }
        \eqsepv
        U_0(s) = \Mcal_{u_0}(s)
        \eqfinp
    \end{equation*}²

Here, $K$ is a convex function, such that $K(1) = 1$. We denote $s_* = \argmin_{s\in \RR} K(s)$. 

The next proposition shows that the behaviour of the repartition of the population depends on $s_* < 1$ or $s_* > 1$. 
\begin{figure}
    \centering
    \includegraphics[scale=0.8]{pictures/severalK.png}
    \caption{Severals K}
    \label{fig:severalK}
\end{figure}

\begin{proposition} Let $R>0$.
    \begin{enumerate}
        \item Suppose that $s_* < 1$, then
            \begin{align*}
            \lim_{t\to \infty} \int_R^\infty u(t,x) \diff x = \int_0^\infty u_0(x) \diff x
            \eqfinp
            \end{align*}
        \item Suppose that $s_* > 1$, then
            \begin{align*}
            \lim_{t\to \infty} \int_0^R u(t,x) \diff x = \int_0^\infty u_0(x) \diff x
            \eqfinp
            \end{align*}
    \end{enumerate}
\end{proposition}

\begin{proof}\text{}\\

\textbf{First case:} Suppose that $s_* < 1$. Let $c$ be a real $>1$. 
By Mellin inversion, we have that
\begin{align*}
    u(t,x) = \frac{1}{2i\pi}\int_{c-i\infty}^{c+i\infty} U_0(s) x^{s} e^{-\lambda t(1 - K(s))} \diff s
\end{align*}

It follows that
\begin{align*}
    \int_R^\infty u(t,x) \diff x 
    &= \frac{1}{2i\pi} \int_R^\infty \int_{c-i\infty}^{c+i\infty} U_0(s) x^{-s} e^{-\lambda t(1 - K(s))} \diff s \diff x \\
    &= \frac{1}{2i\pi} \int_{c-i\infty}^{c+i\infty} \int_R^\infty U_0(s) x^{-s} e^{-\lambda t(1 - K(s))} \diff x \diff s \\
    &= \frac{1}{2i\pi} \int_{c-i\infty}^{c+i\infty} U_0(s) \frac{R^{-s+1}}{s-1} e^{-\lambda t(1 - K(s))} \diff s \\
    &= \frac{1}{2i\pi} \int_{c-i\infty}^{c+i\infty}\frac{U_0(s)}{s-1} e^{-\lambda t(1 - K(s))} e^{(-s+1)\ln R} \diff s 
\end{align*}
(The exchange between the two integrals come from the fact that $c > 1$.)

We deform the integration contour towards the left and cross the pole $s = 1$. The Cauchy's residue theorem ensures that
\begin{align}
    \int_R^\infty u(t,x) \diff x &= \mathrm{Res(1)} + \frac{1}{2i\pi} \int_{s_*-i\infty}^{s_*+i\infty} \frac{U_0(s)}{s-1} e^{-\lambda t(1 - K(s))} e^{(-s+1)\ln R} \diff s
    \label{eq:cauchy1}
\end{align}

We define
\begin{align*}
    \Rcal(t) = \frac{1}{2i\pi} \int_{s_*-i\infty}^{s_*+i\infty} \frac{U_0(s)}{s-1} e^{-\lambda t(1 - K(s))} e^{(-s+1)\ln R} \diff s
    \eqfinp
\end{align*}
We have 
\begin{align*}
    |\Rcal(t)| 
    &\leq \frac{R^{1-s_*}}{2\pi} \int_{s_*-i\infty}^{s_*+i\infty} \left|\frac{U_0(s) }{s-1}\right| \cdot \left|e^{-\lambda t(1 - K(s))}\right| \diff s \\
    &= \frac{R^{1-s_*}}{2\pi} \int_{s_*-i\infty}^{s_*+i\infty} \left|\frac{U_0(s) }{s-1}\right| e^{\Re[-\lambda t(1 - K(s))]} \diff s
    \eqfinp
\end{align*}
However, for all $s = s_* + i\nu \in (s_*-i \infty, s_* + i \infty)$
\begin{align*}
    \Re[-\lambda t(1 - K(s))] 
    &= \Re[-\lambda t(K(1) - K(s))] \\
    &= -\lambda t \Re\left( K(1) - \int_0^\infty z^{s-1}k_0(z)\diff z  \right) \\
    &= -\lambda t \Re\left( K(1) - \int_0^\infty z^{s-1}k_0(z)\diff z  \right) \\
    &=  -\lambda t\left( K(1) - \int_0^\infty \cos(\nu \log z) z^{s_*-1}k_0(z)\diff z  \right) \\
    &\leq -\lambda t \left( K(1) - \int_0^\infty z^{s_*-1}k_0(z)\diff z  \right) \\
    &= -\lambda t \left( K(1) - K(s_*)  \right)
\end{align*}
By Taylor's expansion in $s_*$, there exists a real $\xi \in (s_*, 1)$ such that 
\begin{align*}
    \Re[-\lambda t(1 - K(s))] = - \lambda t \frac{(s_*-1)^2}{2}  K''(\xi) 
\end{align*}
As $ K''(\xi) > 0$, it follows that 
\begin{align*}
    |\Rcal(t)| 
    &\leq \frac{R^{1-s_*}}{2\pi} \int_{s_*-i\infty}^{s_*+i\infty} \left|\frac{U_0(s) }{s-1}\right| e^{\Re[-\lambda t(1 - K(s))]} \diff s
    \eqfinp \\
    &= \frac{R^{1-s_*}}{2\pi} \int_{s_*-i\infty}^{s_*+i\infty} \left|\frac{U_0(s) }{s-1}\right| e^{- \lambda t \frac{(s_*-1)^2}{2}  K''(\xi) } \diff s \\
    &= e^{- \lambda t \frac{(s_*-1)^2}{2}  K''(\xi)} \times \frac{R^{1-s_*}}{2\pi} \int_{s_*-i\infty}^{s_*+i\infty} \left|\frac{U_0(s) }{s-1}\right| \diff s \\
    &\to_{t \to \infty} 0
    \eqfinp
\end{align*}

Using Equation~\ref{eq:cauchy1}, it follows that 

\begin{align*}
    \int_R^\infty u(t,x) \diff x &= \mathrm{Res(1)} + \Rcal(t) \\
    &= \lim_{s \to 1} (s-1)\frac{U_0(s)}{s-1} e^{-\lambda t(1 - K(s))} e^{(-s+1)\ln R}  + \Rcal(t)\\
    &= U_0(1) + \Rcal(t)\\
    &\to_{t \to \infty} \int_0^\infty u_0(x) \diff x
    \eqfinp
\end{align*}

\textbf{Second case:} Suppose that $s_* > 1$. Let $c$ be a real $<1$. \\
By Mellin inversion, we have that
\begin{align*}
    u(t,x) = \frac{1}{2i\pi}\int_{c-i\infty}^{c+i\infty} U_0(s) x^{s} e^{-\lambda t(1 - K(s))} \diff s
\end{align*}

It follows that
\begin{align*}
    \int_0^R u(t,x) \diff x 
    &= \frac{1}{2i\pi} \int_0^R \int_{c-i\infty}^{c+i\infty} U_0(s) x^{-s} e^{-\lambda t(1 - K(s))} \diff s \diff x \\
    &= \frac{1}{2i\pi} \int_{c-i\infty}^{c+i\infty} \int_0^R U_0(s) x^{-s} e^{-\lambda t(1 - K(s))} \diff x \diff s \\
    &= \frac{1}{2i\pi} \int_{c-i\infty}^{c+i\infty} U_0(s) \frac{R^{-s+1}}{1-s} e^{-\lambda t(1 - K(s))} \diff s \\
    &= \frac{1}{2i\pi} \int_{c-i\infty}^{c+i\infty} \frac{U_0(s)}{1-s} e^{-\lambda t(1 - K(s))} e^{(1-s)\ln R} \diff s \\
\end{align*}
We deform the integration contour towards the right and cross the pole $s = 1$. The Cauchy's residue theorem ensures that
\begin{align}
    \int_0^R u(t,x) \diff x &= -\mathrm{Res(1)} + \frac{1}{2i\pi} \int_{s_*-i\infty}^{s_*+i\infty} \frac{U_0(s)}{1-s} e^{-\lambda t(1 - K(s))} e^{(-s+1)\ln R} \diff s
    \label{eq:cauchy2}
\end{align}

As previously, we define
\begin{align*}
    \Rcal(t) = \frac{1}{2i\pi} \int_{s_*-i\infty}^{s_*+i\infty} \frac{U_0(s)}{1-s} e^{-\lambda t(1 - K(s))} e^{(-s+1)\ln R} \diff s
    \eqfinp
\end{align*}
and we can show that $\lim_{t\to\infty}\Rcal(t) = 0$. Using Equation~\ref{eq:cauchy2}, it follows that 

\begin{align*}
    \int_0^R u(t,x) \diff x &= -\mathrm{Res(1)} + \Rcal(t) \\
    &= -\lim_{s \to 1} (s-1)\frac{U_0(s)}{1-s} e^{-\lambda t(1 - K(s))} e^{(-s+1)\ln R}  + \Rcal(t)\\
    &= U_0(1) + \Rcal(t)\\
    &\to_{t \to \infty} \int_0^\infty u_0(x) \diff x
    \eqfinp
\end{align*}
\end{proof}

\subsection{Location of the moment}

\begin{proposition} Let $n>1$ be an integer. Suppose that $K$ is minimum in $s_* > 1$, then for all $R > 0$
    \begin{enumerate}
        \item If $s_* > n$ then 
    \begin{align*}
        \lim_{t\to \infty} \int_0^R x^{n-1} u(t,x) \diff x = 0
        \eqfinp
    \end{align*}
        \item If $s_* < n$ then 
        \begin{enumerate}
            \item If $K(n) < 1$
            \begin{align*}
                \lim_{t\to \infty} \int_R^{\infty} x^{n-1} u(t,x) \diff x = 0
                \eqfinp
            \end{align*}
            \item If $K(n) > 1$
            \begin{align*}
                \lim_{t\to \infty} \int_R^{\infty} x^{n-1} u(t,x) \diff x = \infty
                \eqfinp
            \end{align*}
        \end{enumerate}
    \end{enumerate}
\end{proposition}

\begin{proof}
\textbf{First case:}

Suppose that $s_* > n$. Let $c$ be a real $<n$. \\
    By Mellin inversion, we have that
    \begin{align*}
        u(t,x) = \frac{1}{2i\pi}\int_{c-i\infty}^{c+i\infty} U_0(s) x^{s} e^{-\lambda t(1 - K(s))} \diff s
    \end{align*}

It follows that
\begin{align*}
    \int_0^R x^{n-1}u(t,x) \diff x 
    &= \frac{1}{2i\pi} \int_0^R \int_{c-i\infty}^{c+i\infty} U_0(s) x^{-s+n-1} e^{-\lambda t(1 - K(s))} \diff s \diff x \\
    &= \frac{1}{2i\pi} \int_{c-i\infty}^{c+i\infty} \int_0^R U_0(s) x^{-s+n-1} e^{-\lambda t(1 - K(s))} \diff x \diff s \\
    &= \frac{1}{2i\pi} \int_{c-i\infty}^{c+i\infty} U_0(s) \frac{R^{-s+n}}{n-s} e^{-\lambda t(1 - K(s))} \diff s \\
    &= \frac{1}{2i\pi} \int_{c-i\infty}^{c+i\infty} \frac{U_0(s)}{n-s} e^{-\lambda t(1 - K(s))} e^{(-s+n)\ln R} \diff s
\end{align*}
We deform the integration contour towards the right and cross the pole $s = n$. The Cauchy's residue theorem ensures that
\begin{align}
    \int_0^R x^{n-1}u(t,x) \diff x &= -\mathrm{Res(n)} + \frac{1}{2i\pi} \int_{s_*-i\infty}^{s_*+i\infty} \frac{U_0(s)}{n-s} e^{-\lambda t(1 - K(s))} e^{(-s+n)\ln R} \diff s 
    \label{eq:cauchyn}
\end{align}


As previously, we define
\begin{align*}
    \Rcal(t) = \frac{1}{2i\pi} \int_{s_*-i\infty}^{s_*+i\infty} \frac{U_0(s)}{n-s} e^{-\lambda t(1 - K(s))} e^{(-s+n)\ln R} \diff s
    \eqfinp
\end{align*}
and we can show that $\lim_{t\to\infty}\Rcal(t) = 0$. Using Equation~\ref{eq:cauchyn}, and the fact that $K(n) < 1$ it follows that 


\begin{align*}
    \int_0^R x^{n-1}u(t,x) \diff x &= -\mathrm{Res(n)} + \Rcal(t) \\
    &= -\lim_{s \to n} (s-n)\frac{U_0(s)}{n-s} e^{-\lambda t(1 - K(s))} e^{(-s+n)\ln R}  + \Rcal(t)\\
    &= U_0(n)e^{-\lambda t(1 - K(n))} + \Rcal(t)\\
    &\to_{t \to \infty} 0
    \eqfinp
\end{align*}

\textbf{Second case:}

Suppose that $s_* < n$. Let $c$ be a real $>n$. \\
    By Mellin inversion, we have that
    \begin{align*}
        u(t,x) = \frac{1}{2i\pi}\int_{c-i\infty}^{c+i\infty} U_0(s) x^{s} e^{-\lambda t(1 - K(s))} \diff s
    \end{align*}

It follows that
\begin{align*}
    \int_0^R x^{n-1}u(t,x) \diff x 
    &= \frac{1}{2i\pi} \int_R^{\infty} \int_{c-i\infty}^{c+i\infty} U_0(s) x^{-s+n-1} e^{-\lambda t(1 - K(s))} \diff s \diff x \\
    &= \frac{1}{2i\pi} \int_{c-i\infty}^{c+i\infty} \int_R^{\infty} U_0(s) x^{-s+n-1} e^{-\lambda t(1 - K(s))} \diff x \diff s \\
    &= \frac{1}{2i\pi} \int_{c-i\infty}^{c+i\infty} U_0(s) \frac{R^{-s+n}}{s-n} e^{-\lambda t(1 - K(s))} \diff s \\
    &= \frac{1}{2i\pi} \int_{c-i\infty}^{c+i\infty} \frac{U_0(s)}{s-n} e^{-\lambda t(1 - K(s))} e^{(-s+n)\ln R} \diff s
\end{align*}
We deform the integration contour towards the left and cross the pole $s = n$. The Cauchy's residue theorem ensures that
\begin{align}
    \int_R^{\infty} x^{n-1}u(t,x) \diff x &= \mathrm{Res(n)} + \frac{1}{2i\pi} \int_{s_*-i\infty}^{s_*+i\infty} \frac{U_0(s)}{s-n} e^{-\lambda t(1 - K(s))} e^{(-s+n)\ln R} \diff s 
    \label{eq:cauchyn}
\end{align}


As previously, we define
\begin{align*}
    \Rcal(t) = \frac{1}{2i\pi} \int_{s_*-i\infty}^{s_*+i\infty} \frac{U_0(s)}{s-n} e^{-\lambda t(1 - K(s))} e^{(-s+n)\ln R} \diff s
    \eqfinp
\end{align*}
and we can show that $\lim_{t\to\infty}\Rcal(t) = 0$. Using Equation~\ref{eq:cauchyn}


\begin{align*}
    \int_0^R x^{n-1}u(t,x) \diff x &= \mathrm{Res(n)} + \Rcal(t) \\
    &= \lim_{s \to n} (s-n)\frac{U_0(s)}{s-n} e^{-\lambda t(1 - K(s))} e^{(-s+n)\ln R}  + \Rcal(t)\\
    &= U_0(n)e^{-\lambda t(1 - K(n))} + \Rcal(t)\\
\end{align*}

It follows that 
\begin{equation}
    \int_0^R x^{n-1}u(t,x) \diff x \to_{t \to \infty}
    \begin{cases}
    0 \text{ if $K(n) < 1$}\\
    \infty \text{ if $K(n) > 1$}
    \end{cases}
    \eqfinp
\end{equation}

\end{proof}

\begin{remark}
    We can change the path of integration because : 
\end{remark}
\begin{proof}
    Let $M > 0$. For any $s_* < c$, we have that 
    \begin{align*}
        \bigg| \int_{s_*+iM}^{c+iM} \frac{U_0(s)}{n-s} e^{-\lambda t(1 - K(s))} &e^{(-s+n)\ln R} \diff s \bigg|
        \\
        &\leq 
     \int_{c+iM}^{s_*+iM} \left|\frac{U_0(s)}{n-s}\right| e^{-\lambda t(1 - K(s_*))} e^{\Re((-s+n))\ln R} \diff s \\
     &= 
     \int_{c}^{s_*} \left|\frac{U_0(s+iM)}{n-s-iM}\right| e^{-\lambda t(1 - K(s_*))} e^{(-s+n)\ln R} \diff s.
    \end{align*}
But, for all $s \in (s_*, c)$
\begin{align*}
    \left|\frac{1}{n-s-iM}\right| = \left|\frac{n-s + iM}{(n-s)^2+M^2}\right| \leq \frac{|n-s|}{M^2} + \frac1M
\end{align*}
It follows that $$ \bigg| \int_{s_*+iM}^{c+iM} \frac{U_0(s)}{n-s} e^{-\lambda t(1 - K(s))} e^{(-s+n)\ln R} \diff s \bigg| \to 0$$ when $M \to \infty$.
\end{proof}

\subsection{An explicit solution}
\label{s:explicit}

We assume that     \begin{equation}
    u_0 \geq 0 \mtext{and} \int_0^\infty u_0(x) (1+x) \dx < \infty
    \eqfinp
    \label{H:u0_L1}
\end{equation}

Define 
\begin{equation}
    I(u_0) := \left\{ p \in \RR \eqsepv \int_0^\infty u_0(x) x^{p-1} \dx < \infty \right\}
    \eqfinp
\end{equation}
By Assumption~\eqref{H:u0_L1}, $[1, 2] \subset I(u_0)$. Moreover $I(u_0)$ is a interval of $\RR$.\\
It leads us to define
\begin{subequations}
\begin{align*}
    p_0 &= \inf \{p : p \in I(u_0) \} \in [-\infty, 1] \eqfinv\\
    q_0 &= \sup \{q : q \in I(u_0) \} \in [2, \infty] \eqfinp
\end{align*}    
\end{subequations}

For the simplicity, we assume that 
$ p_0=-\infty$ and $q_0 = \infty$

We also assume that 
\begin{equation}
    \forall \nu \in (p_0, q_0) 
    \eqsepv
    \int_0^\infty |u_0'(x)| x^\nu \dx < \infty
    \eqsepv
    \lim_{x \to 0} x^\nu u_0(x) = 0
    \eqsepv
    \lim_{x \to \infty} x^\nu u_0(x) = 0
    \label{H:DoumicEsco22}
\end{equation}
\begin{equation}
    \forall \nu \in (p_0, q_0) 
    \eqsepv
    \int_0^\infty |u_0''(x)| x^\nu \dx < \infty
    \eqsepv
    \lim_{x \to 0} x^{\nu + 1} u_0'(x) = 0
    \eqsepv
    \lim_{x \to \infty} x^{\nu + 1} u_0'(x) = 0
    \label{H:DoumicEsco23}
\end{equation}

\begin{theorem}(Existence of an explicit solution) Assume \eqref{H:u0_L1}, \eqref{H:DoumicEsco22} and \eqref{H:DoumicEsco23}.
    Then, there exists a unique function $u \in \Ccal([0, \infty); L^1((1+x)dx)) \cap \Ccal^1((0, \infty); L^1((1+x)dx))$ that satisfies Equation~\eqref{eq:main} for all $t > 0$ and almost every $x > 0$. This solution is given by \eqref{eq:explicit_reconstruction} for all $\nu \in (p_0, q_0)$.
    \end{theorem}
    \begin{proof}
    The result is a consequence of the Mellin inversion formula.
    In order to apply it, we have to prove that for any $\nu \in (p_0, q_0)$, 
        \begin{equation*}
            \int\limits_{\nu - i \infty}^{\nu + i \infty} \bigg| U_0(s) e^{\lambda(K(s) - 1) t} x^{-s} \bigg| \ds < \infty
            \eqfinp
        \end{equation*}
    We know that for all $s \in (\nu - i\infty, \nu + i\infty)$, 
        \begin{equation*}
             \bigg| U_0(s) e^{\lambda(K(s) - 1) t} x^{-s} \bigg|  <   \bigg| U_0(s)  e^{\lambda(K(\nu) - 1) t}  x^{-\nu} \bigg|
        \end{equation*}
    
    It only remains to prove that 
        \begin{equation}
                \int\limits_{\nu - i \infty}^{\nu + i \infty} \big| U_0(s)\big| \ds 
                = 
                \int\limits_{-\infty}^{\infty} \big| U_0(\nu + is)\big| \ds
                < \infty
                \label{eq:DouEsc43}
        \end{equation}
    
    This quantity has already been studied under the same assumptions by Doumic and Escobedo \cite{doumic2016time}. Integration by parts lead to 
        \begin{align*}
            U_0(\nu + i v) 
            &= \int_0^\infty u_0(x) x^{\nu-1}x^{i v}\dx \\
            &= \frac{1}{iv + 1} \big[ u_0(x) x^{\nu} x^{i v} \big]_0^\infty - \frac{1}{iv + 1} \int_0^\infty (u_0(x) x^{\nu-1})' x^{i v+1}\dx \\
            \eqfinp
        \end{align*}
    Assumption \eqref{H:DoumicEsco22} guarantees that $[ u_0(x) x^{\nu} x^{i v} ]_0^\infty = 0$. Therefore
        \begin{align*}
            U_0(\nu + i v) 
            &= - \frac{1}{iv + 1} \int_0^\infty (u_0(x) x^{\nu-1})' x^{i v+1}\dx \\
            &= - \frac{\nu - 1}{iv + 1} \int_0^\infty u_0(x) x^{\nu-2} x^{i v + 1}\dx  - \frac{1}{iv + 1} \int_0^\infty u_0'(x) x^{\nu-1} x^{i v + 1}\dx \\
            &= - \frac{\nu - 1}{iv + 1} \int_0^\infty u_0(x) x^{\nu-1} x^{i v}\dx  - \frac{1}{iv + 1} \int_0^\infty u_0'(x) x^{\nu} x^{i v}\dx \\
            &= - \frac{\nu - 1}{iv + 1} U_0(\nu + i v)   - \frac{1}{iv + 1} \int_0^\infty u_0'(x) x^{\nu} x^{i v}\dx
            \eqfinp
        \end{align*}
    It follows that 
        \begin{align*}
            U_0(\nu + i v)
            &= - \frac{1}{(iv + \nu)}\int_0^\infty u_0'(x) x^{\nu} x^{i v}\dx
            \eqfinp
        \end{align*}
    \onlyME{It follow that 
        \begin{equation*}
            |U_0(\nu + i v)| \leq \frac{1}{\sqrt{\nu^2 + v^2}} \int_0^\infty |u_0'(x)| x^{\nu} \dx
            \eqfinp
        \end{equation*}}
    The same calculations and Assumption\eqref{H:DoumicEsco23} allow us to obtain 
        \begin{align*}
            U_0(\nu + i v)
            &=  \frac{1}{(iv + \nu)(iv + \nu + 1)}\int_0^\infty u_0''(x) x^{\nu + 1} x^{i v}\dx
            \eqfinp
        \end{align*}
    It follow that 
        \begin{equation*}
            |U_0(\nu + i v)| \leq \frac{1}{\sqrt{(\nu^2 + v^2)((1+\nu)^2 + v^2)}} \int_0^\infty |u_0''(x)| x^{\nu + 1} \dx
            \eqfinp
        \end{equation*}
    
    We have that $\int_0^\infty |u_0''(x)| x^{\nu + 1} \dx < \infty$ from Assumption\eqref{H:DoumicEsco23}. It allows us to deduce \eqref{eq:DouEsc43} \\
    
    Conversely, if $u_0$ satisfies \eqref{H:u0_L1}, \eqref{H:DoumicEsco22} and \eqref{H:DoumicEsco23}, then 
     $$u \in \Ccal([0, \infty); L^1((1+x)dx)) \cap \Ccal^1((0, \infty); L^1((1+x)dx))$$
    where $u$ is defined by \eqref{eq:explicit_reconstruction}. By construction, it is a solution of the equation \eqref{eq:main}.
    \end{proof}

\section{The behaviour of $u(t,x)$}
We introduce the function: 
    \begin{align*}
        w(t,x) = e^t u(t,x)
        \eqfinp
    \end{align*}
It means that
    \begin{equation*}
        \omega(t,x)
        = \int_{c-i\infty}^{c+i\infty} U_0(s) e^{\lambda K(s)t}x^{-s} \ds
        =  \int_{c-i\infty}^{c+i\infty} U_0(s) e^{\phi(s, t, x)} \ds
        \eqfinp
    \end{equation*}
with $\phi(s,t,x) := \lambda K(s) t - s \ln(x) $.\\

In order to apply the steepest descent method, our first step consist in finding the stationary point of $\phi'(\cdot, t, x) = 0$. 

Define $ s_+(t,x) := (K')^{-1} \Big( \frac{\log x}{\lambda t} \Big)$.


\begin{theorem}
    Suppose that  $u_0$ satisfies the conditions : ...
    Let $\varepsilon(t)$ be any function of $t$ such that $\varepsilon(t) \to 0$ but $t\varepsilon^2(t) \to \infty$ as $t \to \infty$
Then for all $\delta > 0$, there exists a function $\omega_\delta$ such that:
\begin{equation*}
    \begin{split}
    u(t,x) = x^{-s_+(t,x)}e^{(K(-s_+(t,x))-1)t} \times 
    \left( \frac{U_0(s_+(t,x)+\omega_\delta(\varepsilon))}{\sqrt{2\pi tK''(s_+)}} \Theta_1(t)\right)
    \end{split}
\end{equation*}
as $t\to \infty$, uniformly in $x$ and where 
\begin{align*}
    \Theta_1(t) = \left(1 + \Ocal \left( \frac{e^{-\frac{t\varepsilon^2(t)}{2}K''(s_+)}}{\sqrt{t\varepsilon^2(t)K''(s_+)}} \right) \right)
\end{align*}
as $t \to \infty$
\end{theorem}

\begin{proof}
We want to apply the \textit{steepest descent method} \cite{debye1909naherungsformeln} when $t \to \infty$. For this purpose, we split $w(t,x)$ in two parts : 
\begin{align*}
    \omega(t,x)
        &= \int_{s_+-i\infty}^{s_++i\infty} U_0(s)e^{\phi(s, t, x)} \ds \\
        &= \int_{s_+-i\varepsilon}^{s_++i\varepsilon} U_0(s) e^{\phi(s, t, x)} \ds 
        +\int_{\Re(s) = s_+, |Im(s_+)| \geq \varepsilon} U_0(s) e^{\phi(s, t, x)} \ds \\
\end{align*}
We want to show that the "whole mass" comes from the first term. Close to $s_+(t,x)$, we consider the quantity
    \begin{align*}
    \frac{1}{2i \pi} \int_{s_+(t,x) - i\varepsilon}^{s_+(t,x) + i\varepsilon} U_0(s) e^{\phi(s,t,x)} \ds 
    &=
    \frac{1}{2i \pi} U_0(s_+(t,x)) \int_{s_+(t,x) - i\varepsilon}^{s_+(t,x) + i\varepsilon} e^{\phi(s,t,x)} \ds
    \\
    & \qquad \qquad \qquad+ 
    \frac{1}{2i \pi} \int_{s_+(t,x) - i\varepsilon}^{s_+(t,x) + i\varepsilon} (U_0(s) -  U_0(s_+(t,x)))  e^{\phi(s,t,x)} \ds \\
    &= I_1 + I_2
    \end{align*}

\paragraph{\textit{Step 0 -- A first lemma:}}\mbox{}\\
We have
\begin{equation}
    \begin{split}
    K(s) = K(s_+) + (s - s_+(t,x)) K'(s_+(t,x)) + \frac{(s - s_+(t,x))^2}{2} K''(s_+(t,x)) \\+ \frac{(s - s_+(t,x))^3}{6} K'''(c(t,x))
    \end{split}
\end{equation}
where $s < c(t,x) < s_+(t,x)$.
    \begin{align*}
         |K'''(c)|
         &= |\int_0^\infty (\log z)^3 k_0(z) e^{(c-1)\log(z)}\dz|
         \leq \int_0^\infty |(\log z)^3| k_0(z)  e^{\Re (c-1)\log(z)}\dz \\
         &\leq  \boxed{\int_0^\infty |(\log z)^3| k_0(z)  e^{(s_+-1)\log(z)}\dz < \infty}
    \end{align*}
    
Given $\varepsilon > 0$, it follows that for all $t > 0$ and for all $x>0$ such that $|s - s_+| < \varepsilon$, 
    \begin{equation}
        \begin{split}
        \Re K(s) = K(s_+) + \frac{(s - s_+(t,x))^2}{2} K''(s_+(t,x)) + O(\varepsilon^3)
        \eqfinp
        \end{split}
    \end{equation}
and 
    \begin{equation}
        \Re \phi(s, t, x) = \phi(s_+, t, x) + \lambda t \Big(\frac{(s - s_+(t,x))^2}{2} K''(s_+(t,x)) + O(\varepsilon^3)\Big)
        \eqfinp
    \end{equation}
\medskip

\paragraph{\textit{First step -- Study of $I_1$}:}\mbox{}\\
For any $s \in (s_+ - i \infty, s_+ + i \infty)$ such that $|s - s_+| < \varepsilon$.
We know that 
    \begin{equation}
        \Re \phi(s, t, x) = \phi(s_+, t, x) + t \Big(\frac{(s - s_+(t,x))^2}{2} K''(s_+(t,x)) + O(\varepsilon^3)\Big)
        \eqfinp
    \end{equation}

Then 
\begin{align*}
   \int_{s_+(t,x) - i\varepsilon}^{s_+(t,x) + i\varepsilon} e^{\Re \phi(s,t,x)} \ds
    &=
    e^{\phi(s_+, t, x)} \cdot \int_{s_+(t,x) - i\varepsilon}^{s_+(t,x) + i\varepsilon} e^{t \Big(\frac{(s - s_+(t,x))^2}{2} K''(s_+(t,x)) + O(\varepsilon^3)\Big)} \ds \\
    &= 
     i e^{\phi(s_+, t, x)} \cdot \int_{- \varepsilon}^{\varepsilon} e^{t \Big(\frac{(s - s_+(t,x))^2}{2} K''(s_+(t,x)) + O(\varepsilon^3)\Big)} \ds \\
    &= 
     i e^{\phi(s_+, t, x)} \cdot \int_{- \varepsilon}^{\varepsilon} e^{- t \Big(\frac{v^2}{2} K''(s_+(t,x)) + O(\varepsilon^3)\Big)} \ds 
\end{align*}

Now we can see that we are back to the classical method of Laplace. The main idea is to transform the previous integral into a Gauss integral to have
    \begin{equation}
    \int_{- \varepsilon}^{\varepsilon} e^{- t \Big(\frac{v^2}{2} K''(s_+(t,x)) + O(\varepsilon^3)\Big)} \ds 
    \sim_{t \to \infty}
    \sqrt{2\pi} \frac{1}{\sqrt{tK''(s_+)}}
    \end{equation}
    
Indeed, 
\begin{align*}
    \int_{- \varepsilon}^{\varepsilon} e^{- t \Big(\frac{v^2}{2} K''(s_+(t,x)) + O(\varepsilon^3)\Big)} 
    &= 
     \frac{1}{\sqrt{tK''(s_+)}} \int_{- \varepsilon \sqrt{tK''(s_+) + \Ocal(\varepsilon})}^{\varepsilon \sqrt{tK''(s_+) + \Ocal(\varepsilon})} e^{- \frac{y^2}{2}}\dy \\
     \intertext{and}
     \int_{- \varepsilon \sqrt{tK''(s_+) + \Ocal(\varepsilon})}^{\varepsilon \sqrt{tK''(s_+) + \Ocal(\varepsilon})} e^{- \frac{y^2}{2}}\dy &= 
    \sqrt{2\pi} - 2 e^{-\frac{t(\varepsilon^2 K''(s_+) + \Ocal(\varepsilon^3))}{2}} \cdot \Ocal\bigg(\frac{1 }{\varepsilon \sqrt{tK''(s_+) + \Ocal(\varepsilon})}\bigg)
    \eqfinp
\end{align*}
It follows that 
    \begin{align*}
     \frac{1}{2i\pi}   \int_{s_+(t,x) - i\varepsilon}^{s_+(t,x) + i\varepsilon} e^{\Re \phi(s,t,x)} \ds 
     &= 
     \frac{1}{2i\pi} i e^{\phi(s_+, t, x)} \cdot \int_{- \varepsilon}^{\varepsilon} e^{- t \Big(\frac{v^2}{2} K''(s_+(t,x)) + O(\varepsilon^3)\Big)} \ds  \\
     &= 
      \frac{1}{2\pi} e^{\phi(s_+, t, x)} \cdot \frac{1}{\sqrt{tK''(s_+)}} \int_{- \varepsilon \sqrt{tK''(s_+) + \Ocal(\varepsilon})}^{\varepsilon \sqrt{tK''(s_+) + \Ocal(\varepsilon})} e^{- \frac{y^2}{2}}\dy \\
     &=
    \frac{1}{2\pi} e^{\phi(s_+, t, x)} \cdot \frac{1}{\sqrt{tK''(s_+)}} \bigg( 
    \sqrt{2\pi} - 2 e^{-\frac{t(\varepsilon^2 K''(s_+) + \Ocal(\varepsilon^3))}{2}} \cdot \Ocal\bigg(\frac{1 }{\varepsilon \sqrt{tK''(s_+) + \Ocal(\varepsilon})}\bigg) \bigg) \\
    \end{align*}
    
Therefore, 
    \begin{equation}
    I_1 = \frac{U_0(s_+(t,x))}{\sqrt{2\pi}} e^{\phi(s_+, t, x)} \cdot \frac{1}{\sqrt{t \lambda K''(s_+)}} \bigg( 
    1 + \Ocal\bigg(\frac{e^{-\frac{t(\varepsilon^2 \lambda K''(s_+))}{2}} }{\varepsilon \sqrt{\lambda tK''(s_+)}}\bigg) \bigg)\\
    \end{equation}

\medskip
\paragraph{\textit{Second step -- Study of $I_2$}:}\mbox{}\\
We have 
    \begin{equation*}
        I_2 = \frac{1}{2i \pi} \int_{s_+(t,x) - i\varepsilon}^{s_+(t,x) + i\varepsilon} (U_0(s) -  U_0(s_+(t,x)))  e^{\phi(s,t,x)} \ds 
    \end{equation*}
    
Here, we want to go back more or less to the previous integral

    \begin{align*}
         |I_2| 
         &= 
         \frac{1}{2i \pi} \sup_{|v| \leq \varepsilon} \big[U_0(s_+ + iv) -  U_0(s_+) \big] \int_{s_+(t,x) - i\varepsilon}^{s_+(t,x) + i\varepsilon} \big| e^{\phi(s,t,x)} \big| \ds \\
         &=
         \sup_{|v| \leq \varepsilon} \big[U_0(s_+ + iv) -  U_0(s_+) \big] \frac{1}{\sqrt{2\pi}} e^{\phi(s_+, t, x)} \cdot \frac{1}{\sqrt{t \lambda K''(s_+)}} \bigg( 
    1 + \Ocal\bigg(\frac{e^{-\frac{t(\varepsilon^2 \lambda K''(s_+))}{2}} }{\varepsilon \sqrt{\lambda tK''(s_+)}}\bigg) \bigg)
    \end{align*}
    
It only remains to study $  \sup_{|v| \leq \varepsilon} \big[U_0(s_+ + iv) -  U_0(s_+) \big]$ \\

By Definition
\begin{align*}
    \big[U_0(s_+ + iv) -  U_0(s_+) \big]
    &= \Big| \int_0^\infty  z^{s_+ - 1}u_0(z) (z^{iv} - 1) \dz \Big| \\
    &\leq \int_0^{\frac{1}{R}} \Big|  z^{s_+ - 1}u_0(z) (z^{iv} - 1) \Big| \dz
        + \int_{\frac{1}{R}}^R \Big|  z^{s_+ - 1}u_0(z) (z^{iv} - 1) \Big| \dz \\
        &\qquad \qquad \qquad \qquad
        + \int_R^{+\infty} \Big|  z^{s_+ - 1}u_0(z) (z^{iv} - 1) \Big| \dz \\
\end{align*}

If $R \to + \infty$, then the first and the third terms vanish. Therefore, we only need to consider the integral $ \int_{\frac{1}{R}}^R \Big|  z^{s_+ - 1}u_0(z) (z^{iv} - 1) \Big| \dz$. We have
    \begin{align*}
         \int_{\frac{1}{R}}^R \Big|  z^{s_+ - 1}u_0(z) (z^{iv} - 1) \Big| \dz 
         &\leq 
          (\varepsilon \log R)\int_{\frac{1}{R}}^R \Big|  z^{s_+ - 1}u_0(z) \Big| \dz \\
         &\leq 
          (\varepsilon \log R)\int_{0}^{\infty} z^{s_+ - 1}u_0(z)  \dz \\
    \end{align*}

\medskip
\paragraph{\textit{Third step -- A second lemma}:}\mbox{}\\
Outside the neighborhood of $s_+(t,x)$ : $|s - s_+(t,x)| < \varepsilon$. We have
\begin{align*}
    \Re(\phi(s_+ + i \nu, t, x) - \phi(s_+, t, x)) 
    &=
    \Re \bigg( \int_0^\infty x^{s_+ - 1}k_0(x) (x^{i\nu} - 1) \dx \bigg) \\
    &=
    \Re \bigg( \int_0^\infty x^{s_+ - 1}k_0(x) (e^{i\nu \log(x)}- 1) \dx \bigg) \\ 
    &= 
    \int_0^\infty x^{s_+ - 1}k_0(x) [\cos(\nu \log(x))- 1] \dx  \\
    &=: G(v) \leq 0
    \eqfinp
\end{align*}

For any $\nu \in \RR \setminus \{0\}$, we have
\begin{equation}
 \int_0^\infty x^{s_+ - 1}k_0(x) [\cos(\nu \log(x))] \dx < \int_0^\infty x^{s_+ - 1}k_0(x) \dx
 \eqfinp
\end{equation}

It follows that $G(\nu) < 0$, then there exists $\varepsilon > 0$, $R >0$ and $\delta > 0$
such that $\sup_{\varepsilon < \nu < R}G(\nu) < - \delta$.\\

It remains to verify that $\lim_{\nu \to +\infty}G(\nu) \neq 0$. 
\begin{align*}
    G(\nu) 
    &=  \int_0^\infty x^{s_+ - 1}k_0(x) [\cos(\nu \log(x))- 1] \dx \\
    &=  \int_{-\infty}^\infty e^{-s_+ y}k_0(e^{-y})(\cos(-\nu y) - 1 )\dy
    \end{align*}
From Riemann–Lebesgue lemma, we have that 
\begin{equation}
    \lim_{\nu \to +\infty}G(\nu) = - \int_{-\infty}^\infty e^{-s_+ y}k_0(e^{-y})\dy = - \int_{0}^\infty x^{s_+ - 1}k_0(x)\dx < 0
\end{equation}

Therefore
\begin{lemma} For all $\varepsilon > 0$, there exist $\gamma(\varepsilon)$ such
\begin{equation}
    \sup_{|\nu| \geq \varepsilon } \Re(\phi(s_+ + i \nu, t, x)) \leq \Re(\phi(s_+, t, x)) - \gamma(\varepsilon)t
    \eqfinp
\end{equation}
and
\begin{equation}
    \sup_{|\nu| \leq \varepsilon } \Re(\phi(s_+ + i \nu, t, x)) \leq \Re(\phi(s_+, t, x))
    \eqfinp
\end{equation}
\end{lemma}

\medskip
\paragraph{\textit{Fourth step -- Control $\int_{\Re(s) = s_+, |Im(s_+)| \geq \varepsilon} U_0(s) e^{\phi(s, t, x)} \ds$}:}\mbox{}\\
We have that 
\begin{align*}
    \int_{\Re(s) = s_+, |Im(s_+)| \geq \varepsilon} U_0(s) e^{\phi(s, t, x)}
    \ds 
    &\leq  e^{\Re(\phi(s_+, t, x)) - \gamma(\varepsilon)t} \int_{\Re(s) = s_+, |Im(s_+)| \geq \varepsilon} U_0(s) \ds \\
\end{align*}
It follows that 
\begin{align*}
    \int_{\Re(s) = s_+, |Im(s_+)| \geq \varepsilon} U_0(s) e^{\phi(s, t, x)} 
    \ds = \Ocal \left( e^{\Re(\phi(s_+, t, x))} \right)
\end{align*}
\end{proof}


\section{Addition of a term of growth}

Consider the equation
\begin{equation*}\left\{
    \begin{array}{ll}
	\frac{\partial}{\partial t}u(t,x) = - \lambda u(t,x) + \lambda \int_{0}^{\infty} k\big( \frac{x}{y}\big)u(t,y) \frac{\dy}{y} + xu(t,x)
	\eqsepv
	& t \in (0, \infty), x \in \RR_+, \lambda \in (0, \infty) \eqfinv
	\\
	u(0, x) = u_0(x)
	& x \in \RR_+
	\eqfinp
    \end{array}
	\right.
\end{equation*}

By applying the Mellin transform, we have
	\begin{equation*}
		\frac{\partial}{\partial t} U(t,s) = - \lambda U(t,s) + \lambda K(s) U(t,s) + U(t,s+1)
		\eqfinp
	\end{equation*}
It follows that
    \begin{equation}
        \frac{\partial}{\partial t} U(t,s) = - \lambda (1-K(s)) U(t,s) + U(t,s+1)
		\eqfinp
        \label{eq:addgrowth}
    \end{equation}

We consider the Laplace transform of $t \mapsto U(t,s)$
$$ L_s(p) = \int_0^\infty U(t,s)e^{-pt}\diff t$$ 

By applying the Laplace transform to Equation \eqref{eq:addgrowth} we have 
\begin{align*}
    pL_s(p) - U_0(s) = -  \lambda (1-K(s)) L_s(p) + L_{s+1}(p)
\end{align*}
It follows that 
\begin{align*}
    L_{s+1}(p) = - U_0(s) + \big[\lambda (1-K(s)) + p \big]L_s(p)
    \eqfinp
\end{align*}

\begin{remark}
We denote $\alpha(s) = \lambda (1-K(s)) + p$. If $s \in \NN$, and if all quantities are well defined, we have that
\begin{align*}
    L_{s+1}(p) = -U_0(s) -\sum_{n = 1}^{s} \prod_{j = n}^s \alpha(s)U_0(s-1) + L_0(p) \prod_{n = 0}^{s}\alpha(n)
\end{align*}
\end{remark}

\bibliographystyle{plain}
\bibliography{oneforall}


\end{document}
